\begin{lemma}[Situation 7.2, Subcase 2.1 outside the Gamma exception]
Let \(D\) be a \noderef{HierarchySituationTwoCaseTwoData}.  Put
\[
 q=C.\mathit{hierarchy}.\mathit{dividerCount},\qquad
 D_{n,q}=\noderef{HierarchyStoredColorsAt}(C.\mathit{hierarchy},q).
\]
Suppose that the paper's exceptional Gamma condition fails exactly, namely
\[
 \neg\exists\eta\in D_{n,q}\quad
   \eta\in\overline\varphi_n(y_{p-1})
   \quad\hbox{and}\quad \theta\in\Gamma^q_\eta .            \tag{1}
\]
Then a contradiction follows.
\end{lemma}
\begin{proof}
Write \(\mathcal B\) for the supplied hierarchy base.  Put
\(H=C.\mathit{hierarchy}\),
\(A=\noderef{HierarchyStarredStage}(H,q)\),
\(\ell=\noderef{HierarchyFullTailLength}(C)\), and
\[
 B=\noderef{HierarchyFullTailPrefix}(C,\ell-2),\qquad
 B^+=\noderef{HierarchyFullTailPrefix}(C,\ell-1).
\]
The exact prefix definitions, the stage chain in
\noderef{TashkinovHierarchy}, and \(\ell\ge2\) give
\[
 A\preccurlyeq B\preccurlyeq B^+\preccurlyeq T,\qquad
 |B|=|T|-2,\quad |B^+|=|T|-1.                              \tag{2}
\]
Apply \noderef{HierarchySituationTwoCaseTwoExclusions}.  Its three
conclusions are
\[
 \alpha\notin\overline\varphi_n(B),\qquad
 \alpha\notin D_n,\qquad
 \alpha\notin\varphi_n\langle T-A\rangle,                  \tag{3}
\]
where the last assertion quantifies over every exact displayed edge
identity at a position at least \(|A|\).  Apply
\noderef{HierarchyCounterexamplePenultimateElementary} to the six fields
of the supplied \noderef{HierarchyMinimumCounterexampleContext}; it gives
\[
              V(B^+)\text{ is elementary under }\varphi_n. \tag{4}
\]

Invoke \noderef{HierarchySituationTwoReorderedTreesCertified} on the
underlying tail datum of \(D\).  Retain its exact
\noderef{HierarchySituationTwoReorderedTrees} \(R=(T^-,T^*)\), the PE or
exact-seed SE/RE augmentation route for each tree, both actual
\noderef{CorrespondingExtendedTashkinovTree} certificates and maximum
properties, and bases \(\mathcal B^-,\mathcal B^*\) with hierarchies
\(H^-,H^*\).  The certificate says, entry by entry,
\[
 B+(e_p,y_p)=T^-,\qquad
 T^-+(e_{p-1},y_{p-1})=T^*,                                \tag{5}
\]
that \(V(T^-)=V(T)\setminus\{y_{p-1}\}\), that
\(V(T^*)=V(T)\), and that both copied hierarchies have the old core,
closure, divider count, internal stages, and numerical divider positions.

We construct actual \noderef{GoodHierarchy} witnesses \(Q^-\) and
\(Q^*\), rather than treating them as consequences of correspondence.
Transport the old Gamma function across the certified divider-count
equality.  Exact agreement of the raw stored series data, core, and every
internal stage identifies each \noderef{HierarchyStoredColorsAt} and each
stage-missing set.  Hence the two-color, disjoint-pair, refresh,
parallel-base, and minus-closed fields are literally the corresponding
old fields.  Avoidance for an index below \(q\) is also old avoidance,
because both ends of its segment are retained stages.

For avoidance at \(q\), fix \(\zeta\in D_{n,q}\) and
\(c\in\Gamma^q_\zeta\).  If \(\zeta\) has a missing witness in \(B\),
its least position and all earlier edge entries are unchanged in both
reorderings, so old avoidance applies.  Suppose it has no witness in
\(B\).  Deleting \(e_{p-1}\) can remove a constraint but cannot add one;
the only individual edge moved earlier is \(e_p\), whose color is
\(\theta\).  Thus only \(c=\theta\) could fail avoidance.  In that case
\(\zeta\) cannot be missing at \(y_{p-1}\), by the exact negation (1).
Old avoidance of the old last edge \(e_p\), located immediately after
\(y_{p-1}\), then forces the first \(\zeta\)-missing position to be at
most the last position of \(B\), contrary to the present assumption.
For any different stored color, pairwise disjointness excludes \(\theta\)
from its Gamma pair.  This verifies avoidance for every stored color in
both reordered trees and completes all fields of \(Q^-\) and \(Q^*\).

We now make the counting choice in manuscript equation~(11).  Apply
\noderef{NineN} to \(C\) at the exact index \(i=\ell-2<\ell\), using
(4) and the large-set field of the minimum context.  It supplies at least
\(2n+11\) colors that are missing on \(B\), avoid the base-extra missing
set, and are unused on the exact prefix segment \(B-A\).  Enlarging the
used-edge exclusion from \(B-A\) to \(T-A\) adds exactly the two
individual final edges \(e_{p-1}\) and \(e_p\).  Removing the first
edge's color loses at most one color, and removing the second edge's color
loses at most one further color, even if those colors coincide.  Finally
\noderef{TashkinovStoredColorsCardinality} gives \(|D_n|\le n\).
Consequently the remaining set has cardinality at least
\[
                    2n+11-1-1-n=n+9>0.                      \tag{6}
\]
Choose a concrete color \(\beta\) from it.  Thus
\[
 \begin{aligned}
 &\beta\in\overline\varphi_n(B),\qquad
 \beta\notin\overline\varphi_n
      (\mathcal B.\mathit{core}-V(T_n)),\\
 &\beta\notin\varphi_n\langle T-A\rangle,\qquad
 \beta\notin D_n.                                          \tag{7}
 \end{aligned}
\]
The first part of (3) and (7) give \(\alpha\ne\beta\).  Let
\(v_\beta\) be the first displayed vertex missing \(\beta\).  It lies in
\(B\), whereas (3), the missingness of \(\alpha\) at \(y_{p-1}\), and
(4) show that the first \(\alpha\)-missing vertex is exactly
\(y_{p-1}\).

We first prove
\[
                    \beta\notin\overline\varphi_n(y_p).     \tag{8}
\]
If instead \(y_p\) missed \(\beta\), then \(v_\beta,y_p\in V(T^-)\)
would be distinct missing witnesses by (2) and (5).  Together with the
certified corresponding ETT, maximum property, base, hierarchy, and the
explicit good hierarchy \(Q^-\), they form an actual
\noderef{HierarchyCounterexample} \(C^-\).  Its rung and divider count
are the old ones and its last starred stage is exactly \(A\).  Equation
(5) gives full-tail length \(\ell-1\), so the finite-minimum definition
of \noderef{HierarchyPathNumber} gives \(p(C^-)\le\ell-1\).  Since the
tail datum gives \(p(C)=\ell\), the first coordinate of
\noderef{HierarchyCounterexampleMeasure} strictly decreases.
Thus \noderef{HierarchyCounterexampleLess}(C^-,C)\), contradicting the
old lexicographic-minimality field.  This proves (8).

Apply \noderef{ChangeLemma} on the original counterexample with ordered
colors \((\beta,\alpha)\).  Both colors are missing on \(B^+\), at
\(v_\beta\) and \(y_{p-1}\).  Since both lie outside \(D_n\), they are
disjoint from \(D_{n,q}\) and supply the required positive-divider or
zero-divider alternative.  The conditional first-position premise is
active because \(\alpha\) is not missing on \(A\): the first
\(\beta\)-witness lies in \(B\), strictly before \(y_{p-1}\), and (7)
excludes \(\beta\) from every individual tail edge before that position.
The remaining inputs are the six context fields and (4).  Therefore
\[
 K_{\alpha,\beta}(v_\beta)
       =K_{\alpha,\beta}(y_{p-1}).                           \tag{9}
\]

For the second application we must use an actual minimum presentation on
\(T^*\).  Since \(V(T^*)=V(T)\), the old two missing witnesses proving
non-elementarity remain.  With the certified corresponding ETT, maximum
property, base, hierarchy, and \(Q^*\), they form a counterexample
\(C^*\).  Its path number is at most its full-tail length \(\ell\).  A
strict inequality would already contradict minimality.  Otherwise its
first measure coordinate equals that of \(C\); its terminal length and
last ordinary hierarchy stage are both unchanged, so its second
coordinate also agrees.  Hence
\(\mathfrak m(C^*)=\mathfrak m(C)\).

Construct a \noderef{HierarchyMinimumCounterexampleContext} \(\mathcal
X^*\) field by field.  Properness and exact one-uncoloredness are the old
facts about \(\varphi_n\).  Rung agreement transports the complete
lower-rung package.  The common rung and divider count let the old
fewer-divider assertion apply verbatim.  Equality of measures transports
lexicographic minimality: anything smaller than \(C^*\) would be smaller
than \(C\).  Finally, correspondence identifies the terminal ladder tree
and its generating coloring and raw stored data, while rung agreement
identifies the numeral \(2n+11\); this transports the large missing-set
bound.  These are all six fields of the context.

Apply \noderef{HierarchyCounterexamplePenultimateElementary} to
\(C^*\) and \(\mathcal X^*\).  Its penultimate full-tail prefix is
exactly \(T^-\) by (5), so
\[
              V(T^-)\text{ is elementary under }\varphi_n. \tag{10}
\]
In the reordered display the first \(\alpha\)-missing vertex is \(y_p\):
(3) excludes a witness in \(B\), and \(y_p\) is appended before
\(y_{p-1}\).  The first \(\beta\)-missing vertex remains \(v_\beta\in
V(B)\).  Equation (7) excludes \(\beta\) from every exact tail edge in
the new order as well, because (5) reorders the same two individual
edges.  Both colors are missing on the prefix \(T^-\); (8) says precisely
that \(\beta\) is present, not missing, at its new last vertex, while
\(\alpha\) is missing there, so their roles and first positions are
unambiguous.  Using the disjointness from the raw stored set, the explicit
first-position and unused premises just checked, (10), and every field of
\(\mathcal X^*\), a second application of \noderef{ChangeLemma} gives
\[
 K_{\alpha,\beta}(v_\beta)
       =K_{\alpha,\beta}(y_p).                              \tag{11}
\]

Equations (9) and (11) put the three vertices
\(v_\beta,y_{p-1},y_p\) in one global two-color component.  They are
distinct: the first lies in \(B\), while the other two occupy the two
successive positions after \(B\), and the vertex map in
\noderef{TreeSequence} is injective.  At \(v_\beta\), the color \(\beta\)
is missing; at each other vertex, \(\alpha\) is missing.  Properness in
\noderef{IsProperEdgeColoring} and
\noderef{TwoColorIncidenceDegreeAtMostTwo} therefore make all three
vertex-nodes endpoints of the component's spanning path supplied by
\noderef{TwoColorComponentSpanningPath}.  A finite path has at most two
endpoints, a contradiction.
\end{proof}
